\documentclass[11pt,a4paper]{article}

\usepackage[utf8]{inputenc}
\usepackage[T1]{fontenc}
\usepackage{amsmath,amssymb}
\usepackage{graphicx}
\usepackage{hyperref}
\usepackage{booktabs}
\usepackage{float}
\usepackage{algorithm}
\usepackage{algpseudocode}
\usepackage{listings}
\usepackage[margin=2.5cm]{geometry}

\lstset{
    language=Python,
    basicstyle=\ttfamily\small,
    keywordstyle=\bfseries,
    breaklines=true,
    frame=single
}

\title{Project Information Theory\\Compact Disc Digital Audio\\[0.5em]\large Group 24}
\author{
    Debrabandere, Mendel \\
    \and
    D'Hooghe, Achiel \\
    \and
    Loones, Louis \\
    \and
    Ureel, Bram
}
\date{\today}

\begin{document}

\maketitle

\section{Task Division}
% Describe how you divided the tasks among group members and who did what.
Goal for 8 april:
\begin{itemize}
    \item Mendel: 5.3 encoding functions
    \item Achiel: 5.2 write the code \& answer questions
    \item Louis: 5.3 questions + 5.1 questions
    \item Bram: 5.3 decoding functions
\end{itemize}


\begin{table}[H]
\centering
\begin{tabular}{ll}
    \toprule
    \textbf{Member} & \textbf{Tasks} \\
    \midrule
    Name 1 & \\
    Name 2 & \\
    Name 3 & \\
    Name 4 & \\
    \bottomrule
\end{tabular}
\end{table}

\section{Audio CD Subcode}

\subsection{Subcode Datarate}
% Q1: What is the datarate of the subcode in subcode blocks/s? And in bit/s?

\subsection{P Channel Flag Bit}
% Q2: Complete Figure 3 with the state of the channel P flag bit.

\subsection{Track Information}
% Q3: How does the CD player know the number of tracks, starting positions, and total playing time? How is this data protected from errors?

\section{Reed-Solomon Code}

\subsection{Encoding}

\subsubsection{Primitive polynomial}
% Q1a: Show that p(D) = D^8 + D^4 + D^3 + D^2 + 1 is a primitive polynomial with root alpha for GF(2^8).

We need two things: $p(D) = D^8 + D^4 + D^3 + D^2 + 1$ must be irreducible over $\text{GF}(2)$, and its root $\alpha$ must have multiplicative order $2^8 - 1 = 255$.

\paragraph{Irreducibility.}
$p(D)$ has no roots in $\text{GF}(2)$: $p(0) = 1 \neq 0$ and $p(1) = 1+1+1+1+1 = 1 \neq 0$ (mod 2).
Since $\deg p = 8$, any factorisation would include a factor of degree $\leq 4$. Checking division by all irreducible polynomials of degree 2, 3, and 4 over $\text{GF}(2)$ yields a nonzero remainder in every case, so $p(D)$ is irreducible.

\paragraph{Primitivity.}
Let $\alpha$ be a root of $p(D)$, giving us the reduction rule $\alpha^8 = \alpha^4 + \alpha^3 + \alpha^2 + 1$.
Since $255 = 3 \times 5 \times 17$, it suffices to check that $\alpha^d \neq 1$ for $d \in \{85, 51, 15\}$. Computing powers via repeated multiplication and reduction:
\begin{align*}
\alpha^{15} &= \alpha^5 + \alpha^2 + \alpha = 38 \neq 1,\\
\alpha^{51} &= 10 \neq 1,\\
\alpha^{85} &= 214 \neq 1.
\end{align*}
Meanwhile $\alpha^{255} = 1$ because $|\text{GF}(2^8)^*| = 255$. So the order of $\alpha$ is exactly 255, and $p(D)$ is primitive.

\subsubsection{Generator polynomial}
% Q1b: Construct g(x) for the C1 RS code. Give coefficients g_i in GF(2^8). What is m_0?

The $C_1$ code is RS(32,28) over $\text{GF}(2^8)$ with $t = 2$ (IEC~60908, Section~16.3). From the parity check matrix $H_p$ in Figure~11 of the standard and the course notes (slide~145), $m_0 = 1$, so the generator polynomial has roots $\alpha, \alpha^2, \alpha^3, \alpha^4$:
\[
g(x) = (x + \alpha)(x + \alpha^2)(x + \alpha^3)(x + \alpha^4)
\]
(using $-1 = 1$ in characteristic 2). Expanding in pairs:
\begin{align*}
(x+\alpha)(x+\alpha^2) &= x^2 + (\alpha+\alpha^2)x + \alpha^3, \\
(x+\alpha^3)(x+\alpha^4) &= x^2 + (\alpha^3+\alpha^4)x + \alpha^7,
\end{align*}
and multiplying these two quadratics:
\begin{align*}
g(x) &= x^4 + (\alpha+\alpha^2+\alpha^3+\alpha^4)\,x^3
+ (\alpha^3+\alpha^4+\alpha^6+\alpha^7)\,x^2 \\
&\quad + (1+\alpha+\alpha^2+\alpha^5+\alpha^6+\alpha^7)\,x
+ (\alpha^2+\alpha^4+\alpha^5+\alpha^6).
\end{align*}

Writing $g(x) = g_0 + g_1 x + g_2 x^2 + g_3 x^3 + g_4 x^4$, the coefficients are:

\begin{table}[H]
\centering
\begin{tabular}{ccc}
\toprule
 & Polynomial form & Power form \\
\midrule
$g_0$ & $\alpha^2+\alpha^4+\alpha^5+\alpha^6$ & $\alpha^{10}$ \\
$g_1$ & $1+\alpha+\alpha^2+\alpha^5+\alpha^6+\alpha^7$ & $\alpha^{81}$ \\
$g_2$ & $\alpha^3+\alpha^4+\alpha^6+\alpha^7$ & $\alpha^{251}$ \\
$g_3$ & $\alpha+\alpha^2+\alpha^3+\alpha^4$ & $\alpha^{76}$ \\
$g_4$ & $1$ & $\alpha^0$ \\
\bottomrule
\end{tabular}
\end{table}

\subsubsection{Check polynomial}
% Q1c: Give the check polynomial h(x). First and last 5 terms sufficient.

For the full RS(255,251) code, $x^{255} + 1 = g(x) \cdot h(x)$ (since $-1 = 1$ in characteristic~2).
The check polynomial $h(x)$ has degree 251; its roots are all $\alpha^i$ with $i \notin \{1,2,3,4\}$, i.e.\ $\alpha^0, \alpha^5, \alpha^6, \ldots, \alpha^{254}$.

Computing the division gives the first and last five terms:
\[
h(x) = x^{251} + \alpha^{76}x^{250} + \alpha^{165}x^{249} + \alpha^{134}x^{248} + \alpha^{147}x^{247} + \cdots + \alpha^{109}x^3 + \alpha^{145}x^2 + \alpha^{61}x + \alpha^{245}
\]

\subsubsection{Minimum Hamming distance}
% Q1d: What is the minimal Hamming distance? Use the MDS property.

RS codes are MDS codes (maximum distance separable), meaning they meet the Singleton bound with equality:
\[
d_{\min} = n - k + 1 = 32 - 28 + 1 = 5
\]
This allows correction of $t = \lfloor(d_{\min}-1)/2\rfloor = 2$ symbol errors, or detection of up to 4.

\subsubsection{Reed-Solomon vs.\ binary BCH}
% Q1e: Why is RS more appropriate for data protection than binary BCH?

RS codes over $\text{GF}(2^8)$ work on 8-bit symbols rather than individual bits. This is a better fit for CD audio for two reasons.

First, burst errors within a single byte count as only one symbol error for RS. A binary BCH code would treat the same burst as up to 8 separate bit errors, requiring much higher correction capability. A burst corrupting all 8 bits of one byte costs the RS decoder just one of its $t=2$ corrections, while a binary BCH code would need $t \geq 8$.

Second, audio data is stored as bytes, so RS codes in $\text{GF}(2^8)$ match the data format directly. The code rate is also better: RS needs only $2t$ parity symbols (here 4 bytes) to correct $t$ symbol errors, whereas a binary BCH code protecting against the same burst lengths would require many more parity bits.

\subsubsection{Implementation of \texttt{makeGenerator()} and \texttt{encode()}}
% Q1f: Describe your implementation.

\texttt{makeGenerator(m, t, m0)} builds the generator polynomial by starting from $g(x) = 1$ and iteratively multiplying by each factor $(x - \alpha^i)$ for $i = m_0, \ldots, m_0+2t-1$, using \texttt{galois.Poly} arithmetic.

\texttt{encode(msg)} performs systematic encoding row by row. For each information word $m(x)$ of length $l$:
\begin{enumerate}
    \item Place $m(x)$ in the high-order positions of a length-$(l + n - k)$ buffer (equivalent to multiplying by $x^{n-k}$).
    \item Compute the remainder $r(x) = m(x) \cdot x^{n-k} \bmod g(x)$.
    \item Output the codeword $[m(x) \mid r(x)]$: the message followed by the parity symbols.
\end{enumerate}

\subsection{Decoding}

\subsubsection{Bit error correction capability}
% Q2a: How many erroneous bits can C1 always correct? Maximum correctable?

$C_1$ corrects up to $t = 2$ symbol errors, with each symbol being 8 bits.

In the worst case, each erroneous bit lands in a different symbol, so $C_1$ can always correct up to 2 erroneous bits. Three bit errors could fall in three different symbols, which exceeds $t=2$.

In the best case, all bit errors are concentrated in at most 2 symbols. Then up to $2 \times 8 = 16$ erroneous bits can be corrected, since the RS decoder only sees 2 corrupted symbols.

\subsubsection{Bit error rate vs.\ symbol error rate}
% Q2b: Relationship between P_b and P_s for GF(2^8).

With independent bit errors at probability $P_b$, a symbol of $m = 8$ bits is erroneous whenever at least one bit is wrong:
\[
P_s = 1 - (1 - P_b)^8
\]
For small $P_b$ this simplifies to $P_s \approx 8\, P_b$.

\subsubsection{Implementation of \texttt{decode()}}
% Q2c: Describe your decoding algorithm.

The decoder implements the PGZ (Peterson-Gorenstein-Zierler) algorithm with Forney's formula. For each received word $r(x)$:

\begin{enumerate}
    \item Compute syndromes $S_j = r(\alpha^{m_0+j})$ for $j = 0, \ldots, 2t-1$. If all are zero, the word is error-free.

    \item Find the error locator polynomial $\Lambda(x)$ via PGZ: try $\nu = t, t-1, \ldots, 1$ errors, building a $\nu \times \nu$ syndrome matrix each time. The first $\nu$ for which the system is solvable gives $\Lambda(x)$.

    \item Chien search: evaluate $\Lambda(\alpha^{-i})$ for each position $i$ in the codeword to find the error locations. If the number of roots does not equal $\nu$, decoding fails.

    \item Forney's algorithm: compute the error evaluator polynomial $\Omega(x) = S(x) \cdot \Lambda(x) \bmod x^{2t}$ and obtain the error values as
    \[
    e_j = \alpha^{i_j(1-m_0)} \cdot \frac{\Omega(\alpha^{-i_j})}{\Lambda'(\alpha^{-i_j})}
    \]
    where $\Lambda'(x)$ is the formal derivative of $\Lambda(x)$.

    \item Correct the received word and verify by recomputing the syndromes. If they are not all zero, the word is flagged as uncorrectable ($\texttt{nERR} = -1$).
\end{enumerate}

\section{Audio Encoding in CDs}

\subsection{CIRC Structure}
% Q1: Explain the function of every component of the CIRC structure (figures 12 and 13 of the standard). Which error types does C1 protect against? And C2?

\subsection{Implementation of AudioCD Functions}
% Q2: Describe your implementation of the AudioCD class functions.

\subsection{Maximum Correctable Burst}
% Q3: Maximum burst duration (bits) assuming C2 corrects 4 erasures. Translate to scratch width at 1.3 m/s. Compare with simulation results.

\subsection{Linear Interpolation}
% Q4: Describe why interpolation of max 8 subsequent samples works. Relation with sample frequency.

\subsection{Performance Comparison}

\subsubsection{Scratch Errors}
% Q5a: Scratches of 100, 3000, and 10000 databits wide, period 600000 bits.

\subsubsection{Random Bit Errors}
% Q5b: Bit error probability p between 0.05 and 0.001. Plot erasure probability and failed interpolation probability for configs 1, 2, and 3.

\subsection{EFM Modulation}
% Q6: What is EFM and why is it used? Could the EFM demodulator help error correction?

\end{document}
